
\section{拓扑性质}

\trConnected、\trConnectedComponent、（不）可约、\trIrreducibleComponent、\trQuasicompact、\trGenerization、\trSpecialization、\trGenericPoint、\trNoetherian{}\trTopologicalSpace{}与\trClosedPoint{}的定义都已在 \S3.6 中给出。你脑中应当为其中每一个观念都备有图像。

\cref{ex:integral-domain-irreducible} 已表明 $\mathbb{A}^n_k$ 是\trIrreducible{}的（这很容易）。这个论证「在粘合之下表现良好」，由此得到：

\textbf{5.1.A.（容易的练习）}\quad 证明：$\mathbb{P}^n_k$ 是\trIrreducible{}的。

\textbf{5.1.B.（练习）}\quad 练习~3.7.F 已证明：对\trAffineScheme，诸\trIrreducible{}闭子集与诸点之间存在双射（把点 $p$ 送到闭子集 $\overline{\{p\}}$ 的映射是双射）。证明：对一般的\trScheme{}这件事也成立。

\textbf{5.1.C.（容易的练习）}\quad 证明：若概形 $X$ 拥有有限\trCover{} $X=\bigcup_{i=1}^{n}\Spec A_i$（其中诸 $A_i$ 皆\trNoetherian），则 $X$ 是\trNoetherian{}\trTopologicalSpace（\S3.6.14）。（我们很快就会把带有这种\trCover{}的概形称作\emph{\trNoetherian{}概形}，\S5.3.4。）提示：证明一个作为\trNoetherian{}子空间之有限并的\trTopologicalSpace{}自身也是\trNoetherian{}的。

于是 $\mathbb{P}^n_k$ 与 $\mathbb{P}^n_{\mathbb{Z}}$ 都是\trNoetherian{}\trTopologicalSpace：我们是用有限多个\trNoetherian{}环的谱粘合起来造出它们的。

\textbf{5.1.D.（容易的练习）}\quad 证明：概形 $X$ 是\trQuasicompact{}的，当且仅当它能写成有限多个\trAffine\trOpenSubscheme{}之并。（于是对任意\trRing{} $A$，$\mathbb{P}^n_A$ 都\trQuasicompact。）

\textbf{5.1.E.（练习：\trQuasicompact{}概形有\trClosedPoint）}\quad 证明：若 $X$ 是\trQuasicompact{}概形，则每个点的闭包中都含有一个\trClosedPoint。证明：$X$ 的每个非空\trClosedSubset{}都含有 $X$ 的一个\trClosedPoint。特别地，每个非空\trQuasicompact{}概形都有\trClosedPoint。（警告：确有不含任何\trClosedPoint{}的非空概形——例如练习~15.2.P——所以你的论证最好当真用到\trQuasicompactness{}假设！）

这条练习日后常按如下方式使用。设概形的点有某条性质 $P$，且这条性质是「开」的（若点 $p$ 具有性质 $P$，则存在 $p$ 的某个开\trNeighborhood{} $U$，使 $U$ 中所有的点都具有性质 $P$），那么要检验一个\trQuasicompact{}概形是否\emph{所有}的点都具有性质 $P$，只需检验诸\trClosedPoint{}即可。（这一哲学的第一个例子是练习~5.2.E。）这在概形与\trVariety{}的经典理论之间建立了联系——传统\trVariety{}的点正是相应概形的\emph{闭}点（本质上靠的是\trNullstellensatz；参看 \S3.6.9 与练习~5.3.F）。在许多好的情形，\trClosedPoint{}是\trDense{}的（比如对\trVariety{}正是如此；再参看 \S3.6.9 与练习~5.3.F），但在一些基本情形中并非如此（参看\cref{ex:closed-points-dense}(b)）。

\textbf{5.1.1.\ \trQuasiSeparated{}概形。}\quad \trQuasiseparatedness{}是个古怪的观念，但对某些人来说正巧派得上用场。（警告：日后我们会发觉，这其实是\emph{\trMorphism}的性质而非概形的性质，\S8.3.1。）不过大多数人可以对这个观念置之不理，因为他们在现实生活中会遇到的概形全都具有这条性质。称一个\trTopologicalSpace{}是\textbf{\trQuasiSeparated}（quasiseparated）的，若任何两个\trQuasicompact{}\trOpenSubset{}的交仍是\trQuasicompact{}的。（名称中「分离性」的动机在于：它是\emph{\trSeparated}（separated）的弱化版本——对后者，任何两个\emph{\trAffine}开集之交都是\trAffine{}的；参看命题~11.2.13。）

\textbf{5.1.F.（短练习）}\quad 证明：概形是\trQuasiSeparated{}的，当且仅当任何两个\trAffine\trOpenSubscheme{}的交都是有限多个\trAffine\trOpenSubscheme{}之并。

\textbf{5.1.G.（练习）}\quad 证明\trAffineScheme{}都是\trQuasiSeparated{}的。（可能的提示：练习~5.1.F。）

\textbf{5.1.H.（练习）}\quad 证明：概形 $X$ 是\trQuasiSeparated{}的，当且仅当\emph{存在} $X$ 的一个由\trAffine\trOpenSubscheme{}构成的\trCover，其中任何两个之交也被\emph{有限}多个\trAffine\trOpenSubscheme{}覆盖。

我们将会看到，\trQuasiseparatedness{}在定理中（与\trQuasicompactness{}相配合）将是一条有用的假设，而且各种有趣的概形——仿射、局部诺特与分离（分别见练习~5.1.G、练习~5.3.A 与练习~11.2.D）——都是\trQuasiSeparated{}的；这将使我们能把定理陈述得更简洁（例如说「若 $X$ \trQuasicompact{}且\trQuasiSeparated」，而不必说「若 $X$ \trQuasicompact，且这条、那条或另外某条性质成立」）。

「\trQuasicompact{}且\trQuasiSeparated」意味着某件具体的事：

\textbf{5.1.I.（练习）}\quad 证明：概形 $X$ \trQuasicompact{}且\trQuasiSeparated{}，当且仅当 $X$ 能被有限多个\trAffine\trOpenSubscheme{}覆盖，其中任何两个之交也被有限多个\trAffine\trOpenSubscheme{}覆盖。

\textbf{5.1.2.}\ \emph{Qcqs。}\quad 当你在一条定理的假设中见到「\trQuasicompact{}且\trQuasiSeparated」时，应当把这当作一个提示：你将会用到上述解读，而且有限性将以本质的方式被用到。这组假设实在太过常见，以致人们常把它简称为 \textbf{qcqs}（例如见 Qcqs 引理~6.2.9）。

\textbf{5.1.J.（容易的练习）}\quad 证明：一切射影 $A$-\trScheme{}都既\trQuasicompact{}又\trQuasiSeparated。（提示：利用定义中的\trGradedRing{}为\trFinitelyGenerated{}这一事实——那有限个\trGenerator{}会把你引向一组覆盖。）

\textbf{5.1.K.（练习：一个非\trQuasiSeparated{}概形）}\quad 设 $X=\Spec k[x_1,x_2,\dots]$，并令 $U$ 为 $X\setminus[\mathfrak{m}]$，其中 $\mathfrak{m}$ 是\trMaximalIdeal{} $(x_1,x_2,\dots)$。取 $X$ 的两份拷贝，沿 $U$ 粘合起来（「原点加倍了的仿射 $\infty$-空间」；「有限维」版本见例~4.4.5 与\cref{ex:doubled-origin-plane}）。证明：所得的概形不\trQuasiSeparated。提示：这个「\trOpenEmbedding」$U\subset X$ 早在\cref{ex:spec-quasicompact}(b) 中就出现过，是\trAffineScheme{}的非\trQuasicompact{}\trOpenSubset{}的一个例子。

\textbf{5.1.3.}\ \emph{\trDimension{}。}\quad 有一个非常重要的拓扑观念是\emph{\trDimension}。（令人惊叹的是，这竟是一个\emph{拓扑}层面的想法。）不过它虽然在直观上十分基本，却更为困难，所以我们把它推迟到第~12~章。
